Nennen sie die Motivation und Arten von versicherungstechnischen Rückstellungen  


§341e HGB
„Versicherungsunternehmen haben versicherungstechnische Rückstellungen auch insoweit zu bilden, wie dies nach vernünftiger
kaufmännischer Beurteilung notwendig ist, um die dauernde Erfüllbarkeit der Verpflichtungen aus den Versicherungsverträgen
sicherzustellen […]"


1. Beitragsüberträge
Beitragsüberträge sind der Bestandteil der Beiträge, die im Bilanzjahr gebucht wurden. 
D.h. die Beitragsüberträge bilden die Abgrenzung der die Folgejahre betreffenden Beitragsteile.
Die Beitragsüberträge stellen ihren eigenen Charakter nach einen passiven Rechnungsabgrenzungsposten für Beiträge dar.

2. Deckungsrückstellung
Gemäß §341f Abs. 1 HGB sind Deckungsrückstellungen sowohl für LVU als auch für Schaden-Unfall-VU, die VG nach Art der Lebensversicherung
betreiben, zu bilden; z.B. Unfallversicherung mit Prämienrückgewähr, Sachenversicherung gegen Einmalbeitrag (lebenslängliche Hausratversicherung)

3. Rückstellung für noch nicht abgewickelte Versicherungsfälle
Die Schadenrückstellungen sind gemäß §341g HGB anzusetzen für Verpflichtungen, die aus bis zum Ende des Geschäftsjahres eingetretenen, aber
noch nicht abgewickelten Versicherungsfällen resultieren.

4. Rückstellung für erfolgsabhängige und erfolgsunabhängige Beitragsrückerstattung (RfB)
Varianten der Überschussbeteiligung, e.g. Beitragsrückerstattung bei nicht Eintreten des Versicherungsfalls

5. Schwankungsrückstellung und ähnliche Rückstellungen (SchwaRü)
Die Schwankungsrückstellung dient zum Ausgleich der Schwankungen im Schadenverlauf künftiger Jahre (§341h HGB, §29 RechVersV). Trotz
eines homogenen und genügend großen versicherungstechnischen Kollektivs gibt es i.d.R. Schwankungen im Schadenbedarf von Rechnungsperiode
zur Rechnungsperiode.
Sie werden nur in der Schaden- und Unfallversicherung sowie in der Rückversicherung gebildet

6. Sonstige versicherungstechnische Rückstellungen
e.g. Stornorückstellung, Rückstellung für drohende Verluste

Für alle Arten bis auf 5, gilt den 
Bruttobetrag zu nennen davon den Anteil für das in Rückdeckung gegebene Versicherungsgeschäft abzuziehen


Nennen sie einige allgemeine Bewertungsvorschriften für versicherungstechnischen Rückstellungen

Nach §253 Abs. 1 und 2 HGB:
Verbindlichkeiten sind zu ihrem Erfüllungsbetrag anzusetzen
Rückstellungen mit einer Restlaufzeit von mehr als einem Jahr sind abzuzinsen
Der Zinssatz ist entsprechend dem der Restlaufzeit entsprechenden durchschnittlichen Marktzinssatz der vergangenen sieben Geschäftsjahre zu
wählen
Gleiches gilt für Rentenverpflichtungen, für die eine Gegenleistung nicht mehr zu erwarten ist.

Der Aspekt, dass Rückstellungen i.H. des nach kaufmännischer Beurteilung notwendigen Erfüllungsbetrag anzusetzen ist, gilt gem. §341e
Abs. 1 HGB nicht für Versicherungsunternehmen.
Grundsätzlich ist eine Einzelbewertung unter Betrachtung der handelsrechtlichen Möglichkeiten zur Gruppenbewertung vorzunehmen
- Ist eine Einzelbewertung nicht anwendbar, sind Näherungsverfahren gem. § 341e Abs. 3 HGB zulässig

„Ist eine Ermittlung des Wertes der künftigen Verpflichtungen und der künftigen Beiträge nicht möglich, hat die Berechnung auf Grund der aufgezinsten
Einnahmen und Ausgaben der vorangegangenen Geschäftsjahre zu erfolgen (= retrospektive Methode).“
z. B. bei der Fondsgebundenen oder Aktienindexgebundenen Lebensversicherung

Im Allg. gilt für Deckungsrückstellungen:
- Die Verpflichtungen aus dem Lebensversicherungsgeschäft bzw. Krankengeschäft nach Art der LV
-  In Höhe ihres versicherungsmathematisch errechneten Wertes
- Einschließlich bereits zugeteilter Überschussanteile mit Ausnahme der verzinslich angesammelten Überschussanteile
- Nach Abzug des versicherungsmathematisch ermittelten Barwerts der künftigen Beiträge zu bilden (= prospektive Methode)

Sicherstellung der dauernden Erfüllbarkeit der Verpflichtungen aus den Versicherungsverträgen beinhaltet u.a. auch 
Berücksichtigung der aufsichtsrechtlichen Vorschriften bzgl. Rechnungsgrundlagen, 
Rechnungszins und die Zuweisung bestimmter Kapitalerträge zu berücksichtigen.


Wie werden versicherungstechnischen Rückstellungen für Beitragsüberträge bestimmt?

Für eine periodengerechte Erfolgsermittlung, müssen die die Beitragsüberträge bestimmt werden denn sie sind im jeweiligen Geschäftsjahr
nicht anrechenbar, d.h. Sie bilden die Abgrenzung welche die Folgejahre betreffen (deshalb auch abgegrenzten Beitragsteile genannt). 

INSERT IMAGE FROM PAGE 23

Berechnung des Übertragsanteils (ohne Kostenabzug)
• Hauptfällig zum Monat HF
Zahlungsweise (Anzahl der Raten p.a.), ZW = 1,2,4 oder 12

\( \text{BÜ} = \frac{HF - 1 \frac{12}{\text{ZW}}\text{INT}((\text{HF} - 1)\frac{\text{ZW}}{12})}{12} \)
wo \( \text{INT}(\cdot) \) abwärtsrundung notiert (floor function)
Um den Übertragsanteilzu berechnen, ist dieser dann mit der Prämie zu multiplizieren.

Eine Unfallversicherung eines Versicherungsunternehmens ist für einen Jahresbeitrag i.H.v. 1.200 € abzuschließen.
i. Ein Versicherungsnehmer schließt eine Unfallversicherung zum 1.November ab und zahlt jedes Quartal einen Teil des Beitrags.
\( \text{HF} = 1.11 = 11; \text{ZW} = 4 → \frac{(11-1-\frac{12}{4}*\text{INT}([11-1]*\frac{4}{12}))}{12}=\frac{1}{12} \implies\) Übertragsanteil \( = 1/12*1200 = 100 \).
i. Ein Versicherungsnehmer schließt eine Unfallversicherung zum 1. Juli ab und zahlt gleich zu Beginn des Vertragsabschlusses einmalig die
Versicherungsprämie.
\(\text{HF} =1.7=7; \text{ZW}=1 \implies \frac{(7-1-\frac{12}{1}*\text{INT}([7-1]*\frac{1}{12}))}{12}=1/2 \implies\) Übertragsanteil \(= \frac{1}{2}\cdot1200=600\).


Was sind versicherungstechnische Deckungsrückstellungen?


Versicherungstechnische Rückstellungen weisen in der Bilanz von Versicherern deren Verpflichtungen aus Versicherungsverträgen entsprechend 
der handelsrechtlichen Bewertung aus. 

Sie müssen im Jahresabschluss auch insoweit gebildet werden, 
wie dies zur Sicherstellung der dauernden Erfüllbarkeit der Verpflichtungen aus den Versicherungsverträgen erforderlich ist. 

Aufgrund dieser Rückstellungen bestimmt sich aufsichtsrechtlich, in welchem Umfang Kapitalanlagen zur Sicherung 
der Ansprüche der Versicherungsnehmer im Insolvenzfall gehalten werden müssen (Sicherungsvermögen) 
und die Eigenmittel, die der Versicherer haben muss. 
Versicherungstechnische Rückstellungen bilden bei Versicherern den überwiegenden Teil der Passivseite der Bilanz.

Beispiele:
- Bei einer Schadenversicherung der Erwartungswert des Schadens.
- Bei einer Leibrente der Erwartunsgwert des Barwerts (diskontiert) der zu zahlenden Rente.
- Bei Unfallversicherung mit Prämienrückgewähr werden alle eingezahlten Beiträge ohne Rücksicht auf etwa bezahlte Unfallentschädigungen bei
  Erreichen eines bestimmten Alters oder bei vorzeitigem Tod des VN zuzüglich einer Überschussbeteiligung zurückerstattet.

Beispiele:
Gemäß §341f Abs. 1 HGB sind Deckungsrückstellungen sowohl für LVU als auch für Schaden-Unfall-VU, die VG nach Art der Lebensversicherung
betreiben, zu bilden; z.B. Unfallversicherung mit Prämienrückgewähr, Sachenversicherung gegen Einmalbeitrag (lebenslängliche
Hausratversicherung)
Bei Unfallversicherung mit Prämienrückgewähr werden alle eingezahlten Beiträge ohne Rücksicht auf etwa bezahlte Unfallentschädigungen bei
Erreichen eines bestimmten Alters oder bei vorzeitigem Tod des VN zuzüglich einer Überschussbeteiligung zurückerstattet werden. Diese
Versicherung stellt eine Kombination aus UnfallV- und LV-Elementen dar, wobei der LV-Teil in vt. Hinsicht als kapitalbildende LV auf den Todes- und
Erlebensfall mit gleichmäßig steigender Versicherungssumme interpretiert werden kann. Für den hier integrierten Sparprozess ist eine DR zu bilden


Was sind versicherungstechnischen Rückstellungen für noch nicht abgewickelte Versicherungsfälle

Schadenrückstellungen dienen der Erfassung von dem Grund und/oder der Höhe nach ungewissen Verbindlichkeiten ggü. VN bzw. ggü.
Geschädigten Dritten aus realisierten wirtschaftlichen Risiken, die in Versicherungsverträgen von VU übernommen worden sind

Die Schadenrückstellungen sind gemäß §341g HGB anzusetzen für Verpflichtungen, die aus bis zum Ende des Geschäftsjahres eingetretenen, aber
noch nicht abgewickelten Versicherungsfällen resultieren.

E.g. Reparaturkostenversicherung, ein Auto muss repariert werden, jedoch braucht die Werkstatt zur Reparatur min. 1 Monat und kann erst dann eine Rechnung austellen.

Zusammensetzung der bilanzierten Schadenrückstellung des saG

- Teil-Schadenrückstellung für am Abschlussstichtag bekannte Versicherungsfälle
  (ohne Renten-VF und bekannte Spätschäden)
- Teil-Schadenrückstellung für am Abschlussstichtag bekannte Renten-
  Versicherungsfälle (Renten-DR)
- Teil-Schadenrückstellung für bekannte und unbekannte Spätschäden
- Teil-Schadenrückstellung für Schadenregulierungsaufwendungen (SchaReKos)


Was sind versicherungstechnischen Schwankungsrückstellung und ähnliche?

Die Schwankungsrückstellung dient zum Ausgleich der Schwankungen im Schadenverlauf künftiger Jahre (§341h HGB, §29 RechVersV). Trotz
eines homogenen und genügend großen versicherungstechnischen Kollektivs gibt es i.d.R. Schwankungen im Schadenbedarf von Rechnungsperiode
zur Rechnungsperiode. Sie werden daher als Instrument zum Risikoausgleich in der Zeit angesehen.
Sie können insbesonder auch sogenannte „Großrisikenrückstellungen“ sein, z.B. für eine
Versicherung von Atomanlagen, Pharma-Risiken, Terrorrisiken.

Die Schwankungsrückstellung wird nur in der Schaden- und Unfallversicherung sowie in der Rückversicherung gebildet. Dabei wird diese
separate für jeden einzelnen größeren Versicherungszweig gebildet (siehe Anlage zu §29 der RechVersV).

Es müssen für jedes Kollektiv, für das eine SchwaRü gebildet wird, so viele GJ in die Betrachtung einbezogen werden, dass in diesem
Ausgleichszeitraum mit angemessener statistischer Sicherheit angenommen werden kann, dass die Summe der Unterschäden gleich der Summe
der Überschäden ist.
Unterschäden sind wenn die Schadenquote des Bilanzjahres die durchschnittliche Schadenquote unterschreitet und Überschäden analog.

Gemäß FB 1 RechVersV ist die SchwaRü und ähnliche Rückstellungen auf der Passivseite unter dem Posten E V auszuweisen
In der GuV: „Veränderung der Schwankungsrückstellung und ähnlicher Rückstellung“

SchRü müssen nur gebildet werden, wenn 2 aus 3 der folgenden Merkmale erfüllt sind
- die Verdienten Beiträge im Durchschnitt der letzten drei Geschäftsjahre (inkl. Bilanzjahr)
  125.000€ übersteigen.
- die Standardabweichung der SQ des
  BOZ von der durchschnittlichen SQ
- die Summe aus SQ und KQ mind. einmal im BOZ 100\% überschritten hat. mind. 5 Prozentpunkte beträgt.


Was ist die Beitragsbestimmung mittels des Äquivalenzprinzips

Barwert aller rechnungsmäßigen Einnahmen (Beiträge) = Barwert aller rechnungsmäßigen Ausgaben (Versicherungsleistungen und Kosten)


Was ist der Risiko-und Sparanteil der Prämie

Alle Beitragsanteile, die kalkulatorisch im aktuellen Jahr nicht für Kosten und Risikobeitrag benötigt werden, machen den sogenannten Sparanteil
(auch Sparbeitrag) aus. Der Sparanteil wird für die Finanzierung zukünftiger Versicherungsleistungen (Garantieverzinsung und
Überschussbeteiligung) benötigt (Rückstellung).

Der Risikobeitrag ist der erwartete, jährliche Aufwand für den Versicherungsfall. Für eine Todesfallleistung bspw. entspricht er i.W. der
Todesfallwahrscheinlichkeit multipliziert mit der zusätzlich zur Deckungsrückstellung zu zahlenden Todesfallsumme.
Der Sparbeitrag kann im Verlauf der Versicherung auch negativ sein. Dies entspricht einer Entnahme aus der Deckungsrückstellung.

Was sind Ursachen für Überschüsse

Wichtigste Gewinnquellen sind:
• Zinsergebnis
• Risikoergebnis
• Kostenergebnis
Weitere Gewinnquellen sind:
• Ergebnis aus dem vorzeitigem Abgang
• Rückversicherungsergebnis
• Übriges Kapitalanlageergebnis


Gehen sie auf das Zinsergebnis als Überschussentstehung ein

Zinsergebnis
• Ist definiert als Differenz zwischen den tatsächlich aus den Kapitalanlagen erzielten (laufenden) Nettoerträgen und dem
kalkulierten Rechnungszins („rechnungsmäßige Zinsen“)
• Ist in der Regel die wichtigste Gewinnquelle
• Die rechnungsmäßigen Zinsen beinhalten Zinsen auf die DR, das Guthaben aus der verzinslichen Ansammlung und auch
Zinsen auf die Pensionsrückstellung
• Ratenzuschläge für Zinsausfall werden den Nettoerträgen zugeschlagen



Gehen sie auf das Risikoergebnis als Überschussentstehung ein

Ist definiert als Differenz zwischen den Versicherungsleistungen, die sich auf Grund der rechnungsmäßigen Annahmen ergäben
und den tatsächlich erbrachten Versicherungsleistungen.
Gründe für Risikogewinne:
• Biometrische Rechnungsgrundlagen (e.g. Sterbetafeln) enthalten Sicherheitszuschläge
• Statistische Schwankungen des Schadenaufwands
• Selektionseffekte durch Gesundheitsprüfung
• Bei Versicherungen mit Todesfallcharakter wirkt sich die höhere Lebenserwartung positiv aus
(bei Erlebensfallcharakter umgekehrt)


Gehen sie auf das Kostenergebnis als Überschussentstehung ein

Ist definiert als Differenz zwischen den aus den Beiträgen zur Verfügung stehenden Anteilen zur Kostendeckung und den
tatsächlichen Aufwendungen des LVU
Tatsächliche Aufwendungen sind:
• Abschlussaufwendungen (AK-Ergebnis meist negativ)
• Verwaltungsaufwendungen (VK-Ergebnis meist positiv)
• Regulierungsaufwendungen (im VK-Ergebnis enthalten)
Insgesamt sollte sich ein positives Kostenergebnis ergeben


Woraus besteht der Rohüberschuss eins Bilanzperiode
Der Rohüberschuss ist das Ergebnis (= Erträge - Aufwand) eines Geschäftsjahres vor Überschussbeteiligung und Beitragsrückerstattung
besteht aus:
• Risikoergebnis und Ergebnis aus dem vorzeitigem Abgang
• Kapitalanlageergebnis
• Kostenergebnis
• Stornoergebnis
• Rückversicherungsergebnis
• Sonstiges Ergebnis

Der Rohüberschuss eines Geschäftsjahres wird verteilt auf
• Direktgutschrift
• Zuführung zur Rückstellung für Beitragsrückerstattung
• Jahresüberschuss

Was sind die Möglichkeiten für die Überschussverwendung beim VU und VN

Bei Überschussbeteiligung bestimmt der Vorstand mit Zustimmung des Aufsichtsrates die für die Überschussbeteiligung der VN
zurückzustellenden Beträge.
Die für die Überschussbeteiligung der VN bestimmten Beträge sind, soweit sie nicht den VN unmittelbar zugeteilt wurden, in eine
Rückstellung für Beitragsrückerstattung (RfB) einzustellen

Innerhalb des VU und seinen Anlegern nimmt die Ausschüttung folgende Formen an

INSERT IMAGE PAGE 56


Zwischen VU und VN kann die Ausschüttung folgende Formen an

INSERT IMAGE PAGE 57


  

Geben Sie Beispiele wie versicherungstechnische Rückstellungen in Bilanzen dargestellt werden

Beispiel Schaden/Unfall-Versicherung

INSERT IMAGE FROM 83


Beispiel Lebensversicherung:


INSERT IMAGE FROM 86



Geben Sie Beispiele wie versicherungstechnische Rückstellungen in GuV dargestellt werden

Beispiel Schaden/Unfall-Versicherung: GuV


INSERT IMAGE FROM 82


Beispiel Lebensversicherung: GuV


INSERT IMAGE FROM 85
